\section{Diseño y ejemplos de tareas de razonamiento}
\subsection{Ejemplo de razonamiento en varios pasos}
Para que el pipeline de self-rewarding aprenda razonamiento en varios pasos, Weston propone dividir el Chain-of-Thought en tres pasos: «decompose → verify → finalize». Por ejemplo:
\begin{itemize}
    \item Prompt: «Analiza un problema algebraico de 4 pasos, mostrando cada paso con variables intermedias».
    \item Draft: enumera la cadena de razonamiento paso a paso y da una respuesta inicial.
    \item Verification: el modelo vuelve a revisar el draft y, cuando no puede verificarlo, genera una idea alternativa.
\end{itemize}
\textit{``You can make it kind of cross check itself and see that it was wrong and write a final verified response.''}(SRT 1188-1190)Este proceso es la aplicación típica de multi-step reasoning + meta judge.

\begin{table}[H]
\centering
\begin{tabular}{p{4cm}p{4cm}p{6cm}}
\toprule
Tipo de tarea & Prompt de ejemplo & Métrica de verificación \\
\midrule
Razonamiento algebraico & «Descompón x+2y=7, x-y=1» & Integridad del Chain-of-Thought, consistencia de la respuesta final \\
Análisis causal & «Describe cómo un agente de RL reflexiona sobre el reward hacking» & Si señala el punto de desviación en los pasos de la operación \\
Juicio estadístico & «Compara la precisión y la confianza de dos modelos» & Consistencia del meta judge al determinar el winner \\
\bottomrule
\end{tabular}
\caption{Tareas típicas de multi-step reasoning en Self-Instruct.}
\end{table}

\begin{knowledgebox}{Control de calidad del razonamiento en varios pasos}
Cada cadena de razonamiento debe llevar consigo una verified response, de modo que el modelo de reward pueda tratar los errores como señal de supervisión, en lugar de premiar solo la extensión o la fluidez. Tanto una reasoning chain demasiado corta como una chain con overthinking deben perder puntos en el cross-check del judge.
\end{knowledgebox}

\subsection{Ejemplo de programación y verificación}
En las tareas de programación, el scene de Self-Instruct suele exigir que el modelo entregue código completo, explicación y pruebas. El pipeline de self-rewarding hace que el modelo, después de generar el código, vuelva a «ejecutar» su propio razonamiento (por ejemplo, simulando casos de prueba) y luego dé una verified response que explique por qué ese código cumple los requisitos. Esa verified response de paso le enseña al modelo a hacer self-evaluate de su output.

\begin{itemize}
    \item Prompt: «Escribe una función de Python que reciba una lista y devuelva los elementos únicos».
    \item Draft: devuelve la función más una breve explicación.
    \item Meta judge: usa el SelfT evaluator para comparar el draft con el verified, y determina cuál explicación es más clara y cuál código es más seguro.
\end{itemize}

\subsection{Escenario simulado del meta judge}
En el deployment real, el meta judge desempeña el papel del «último guardrail». Por ejemplo, una vez que el draft y el verified están terminados, el judge calcula la diferencia de alignment entre ambos y da un tick con alta confianza. Este check a veces requiere introducir un external agent (como un evaluador humano o de terceros), sobre todo en los casos en que el judge tiene confianza pero hallucinates (SRT 2500-2509).

\subsection{Resumen de la sección}
Mediante multi-step reasoning, tareas de programación y escenarios simulados de meta judge, se puede tratar el entrenamiento de self-rewarding como un diseño de alto nivel, que va encadenando prompt → draft → verified → judge hasta formar una plantilla de entrenamiento repetible.

